\section{Thiết kế gateway ESP32 và ThingsBoard}
Gateway ESP32 là cầu nối giữa mạng LoRa cục bộ và Internet. ESP32 giao tiếp với module LoRa qua SPI hoặc UART tùy module sử dụng. Khi nhận gói tin từ node, gateway tách chuỗi dữ liệu, kiểm tra số trường, chuyển đổi kiểu dữ liệu và tạo payload gửi lên ThingsBoard.

\subsection{Kết nối ThingsBoard}
Gateway được cấu hình token thiết bị của ThingsBoard. Dữ liệu telemetry được gửi theo dạng JSON, trong đó mỗi khóa biểu diễn một thông số đo như \texttt{temperature}, \texttt{humidity}, \texttt{co2}, \texttt{lux}, \texttt{tds} và \texttt{ph}.
Trên ThingsBoard, dashboard hiển thị các widget gồm đồng hồ đo, biểu đồ thời gian, trạng thái quạt/đèn và cảnh báo khi thông số vượt ngưỡng. Các ngưỡng có thể được đặt theo nhu cầu sử dụng, chẳng hạn CO2 cao, ánh sáng thấp hoặc pH nằm ngoài khoảng phù hợp cho cây.

\subsection{Mô hình thiết bị trên ThingsBoard}
Trong ThingsBoard, gateway hoặc node được đại diện bởi một thiết bị có access token riêng. Telemetry được gửi lên theo từng khóa. Cách đặt tên khóa cần thống nhất với app Flutter để tránh lỗi khi truy vấn dữ liệu. Các khóa đề xuất gồm \texttt{temperature}, \texttt{humidity}, \texttt{co2}, \texttt{lux}, \texttt{tds}, \texttt{ph}, \texttt{rl1}, \texttt{rl2}, \texttt{acPower} và \texttt{acSetpoint}.

\begin{longtable}{|p{3.5cm}|p{3cm}|p{6cm}|}
\hline
\textbf{Telemetry} & \textbf{Đơn vị} & \textbf{Ý nghĩa} \\
\hline
\texttt{temperature} & $^\circ$C & Nhiệt độ không khí tại khu vực node. \\
\hline
\texttt{humidity} & \% & Độ ẩm tương đối của không khí. \\
\hline
\texttt{co2} & ppm & Nồng độ CO2, dùng đánh giá mức thông gió. \\
\hline
\texttt{lux} & lux & Cường độ ánh sáng tại vị trí cây. \\
\hline
\texttt{tds} & ppm & Tổng chất rắn hòa tan trong nước hoặc dung dịch. \\
\hline
\texttt{ph} & không đơn vị & Độ pH của dung dịch chăm sóc cây. \\
\hline
\texttt{rl1}, \texttt{rl2} & ON/OFF & Trạng thái relay. \\
\hline
\texttt{acPower} & ON/OFF & Trạng thái dự đoán của điều hòa. \\
\hline
\end{longtable}

\subsection{Rule chain và cảnh báo}
ThingsBoard có thể xử lý dữ liệu thông qua rule chain. Khi telemetry mới được gửi lên, rule chain kiểm tra ngưỡng và tạo cảnh báo nếu cần. Ví dụ, nếu CO2 vượt ngưỡng, hệ thống có thể tạo alarm "CO2 cao"; nếu ánh sáng thấp, hệ thống tạo alarm "Thiếu sáng"; nếu pH nằm ngoài khoảng cho phép, hệ thống tạo alarm "pH bất thường".

Các cảnh báo này giúp người dùng không cần quan sát app liên tục. Trong phiên bản mở rộng, ThingsBoard có thể gửi notification hoặc app Flutter có thể lấy danh sách alarm để hiển thị nổi bật. Đây là điểm quan trọng vì hệ thống kiểm soát chất lượng không khí cần phản ứng với tình huống bất thường, không chỉ hiển thị dữ liệu.

\subsection{Luồng điều khiển xuống node}
ThingsBoard cung cấp cơ chế RPC hoặc shared attributes để gửi lệnh xuống gateway. Gateway lắng nghe lệnh, chuyển thành gói LoRa dạng:
\begin{center}
\texttt{CMD,NODE\_ID,DEVICE,ACTION}
\end{center}
Ví dụ \texttt{CMD,1,FAN,ON} để bật quạt, \texttt{CMD,1,LIGHT,OFF} để tắt đèn, \texttt{CMD,1,AC,TEMP\_UP} để tăng nhiệt độ điều hòa và \texttt{CMD,1,AC,TEMP\_DOWN} để giảm nhiệt độ điều hòa. Node nhận lệnh, thực thi và phản hồi trạng thái.

\subsection{Xử lý mất kết nối}
Gateway cần xử lý hai loại mất kết nối: mất LoRa với node và mất Internet với ThingsBoard. Khi mất LoRa, gateway không nhận được dữ liệu mới và có thể đánh dấu node offline sau một khoảng timeout. Khi mất Internet, gateway vẫn có thể nhận dữ liệu LoRa nhưng không gửi được lên ThingsBoard. Trong trường hợp này, gateway nên có bộ đệm cục bộ để lưu tạm dữ liệu.

\begin{lstlisting}[caption={Mã giả xử lý gửi telemetry tại gateway}]
if (wifiConnected() && thingsboardConnected()) {
    sendTelemetry(payload);
} else {
    cacheTelemetry(payload);
}

if (connectionRecovered()) {
    flushCachedTelemetry();
}
\end{lstlisting}
